\clearpage\thispagestyle{empty}
\pdfbookmark[0]{Sources for Comet (17 May 1821)}{cometmay1821-sources}
\begin{center}\textbf{Sources for Comet (17 May 1821)}\end{center}
\noindent This AI translation follows the report \textit{Comet.} in \href{https://gdz.sub.uni-goettingen.de/id/PPN236007467}{\textit{Werke}, vol.~VI} (Göttingen, 1874), pp.~437--439. The first printing appeared under Göttingen in \textit{Göttingische gelehrte Anzeigen}, no.~78, 17 May 1821, pp.~\href{https://images.p.sub.uni-goettingen.de/iiif/image/gdz:PPN319721507_1821:00000773/full/1600,/0/default.jpg}{769}, \href{https://images.p.sub.uni-goettingen.de/iiif/image/gdz:PPN319721507_1821:00000774/full/1600,/0/default.jpg}{770} and \href{https://images.p.sub.uni-goettingen.de/iiif/image/gdz:PPN319721507_1821:00000775/full/1600,/0/default.jpg}{771}. The digitized copy is from the Göttingen State and University Library. The report presents the Göttingen observations and credits von Staudt with the parabolic elements and brightness calculation.
\par\smallskip\noindent Both printings put the Göttingen meridian on the row for the logarithm of the least distance. This placement and the 10 February declination follow the collected text. In the first observation, the first printing uses a raised u and prime marks where the collected edition has raised u, m and s. The first printing uses decimal commas. The \hyperlink{native-440-part-lower}{report of 10 February 1821} is an earlier Göttingen account of the same comet. Translation and typesetting produced by AI.
\clearpage
